\section{Autopilot Description}

Two autopilots were designed with the same architecture used in the course script.

\subsection{Attitude Autopilot}

The attitude autopilot uses an inner pitch-rate loop and an outer attitude loop. The inner loop feeds back \(q\), measured by the gyro, and acts through the actuator and the short-period dynamics. The outer loop compares the attitude command with \(\theta\), and the plant from \(q\) to \(\theta\) includes the kinematic integrator \(1/s\).

\[
    \theta_c \rightarrow K_\theta \rightarrow
    \left(K_q,~\text{actuator},~q/\delta\right)
    \rightarrow \frac{1}{s} \rightarrow \theta
\]

The inner gain \(K_q\) was selected from the root locus of \(q/\delta\). The outer gain \(K_\theta\) was selected from the attitude-loop root locus, using the damping criterion from the course code.

\subsection{Acceleration Autopilot}

The acceleration autopilot keeps the same inner \(q\)-feedback loop. The external loop controls lateral acceleration \(A_z\), with an integrator in the external amplifier, as used in the course material:

\[
    A_{z,c} \rightarrow \frac{K_a}{s} \rightarrow
    \left(K_q,~\text{actuator},~q/\delta\right)
    \rightarrow A_z/\delta \rightarrow A_z
\]

The first value of \(K_a\) was obtained from the acceleration root locus. A final adjustment was then made around that value to keep the acceleration loop stable while satisfying the minimum frequency-response requirement.

\subsection{Actuator and Airframe Data}

The actuator was kept as the second-order model from \texttt{ProjRLocus.m}, with natural frequency \(100~\text{Hz}\) and damping \(\sqrt{2}/2\). The lateral acceleration limit used in the missile data file is \(5g\), and the fin deflection limit used by \texttt{DadosControle.m} is \(10^\circ\).
